\chapter{Primeiros passos}
\label{chap:primeiros-passos}

\section{O que é LaTeX?}

LaTeX é um sistema de preparação de documentos. Em vez de posicionar cada
elemento visualmente, você escreve texto e comandos que descrevem sua função:
isto é um título, aquilo é uma seção, esta frase merece ênfase. O mecanismo de
compilação interpreta o arquivo-fonte e produz um PDF consistente.

Essa separação entre conteúdo e apresentação é especialmente útil em livros,
artigos, trabalhos acadêmicos e documentos técnicos. Num projeto longo, mudar a
aparência de todos os capítulos deixa de ser uma tarefa manual: altera-se uma
regra compartilhada. LaTeX\index{LaTeX} também automatiza sumário, numeração, referências
cruzadas, bibliografia e índice remissivo \citep{lamport1994}.

Os arquivos de texto normalmente terminam em \arquivo{.tex}. Eles podem ser
editados no VS Code, TeXstudio, Overleaf ou qualquer editor de texto simples.
Uma distribuição como TeX Live fornece o compilador e os pacotes.

\section{Seu primeiro documento}

Crie uma pasta chamada \arquivo{meu-livro} e, dentro dela, um arquivo
\arquivo{main.tex} com este conteúdo:

\begin{lstlisting}[language=LaTeXBook]
\documentclass{article}

\begin{document}
Ola, LaTeX!
\end{document}
\end{lstlisting}

A declaração \comando{documentclass\{article\}} escolhe a classe do documento.
Ela define decisões básicas de layout e comandos disponíveis. O ambiente
\arquivo{document} delimita o conteúdo que aparecerá no PDF. Tudo antes dele é
o \emph{preâmbulo}, onde configuramos pacotes e opções.

Compile no terminal:

\begin{lstlisting}
latexmk -lualatex main.tex
\end{lstlisting}

O arquivo \arquivo{main.pdf} deve surgir na mesma pasta. O comando
\arquivo{latexmk} executa quantas passagens forem necessárias. A opção
\arquivo{-lualatex} seleciona LuaLaTeX, mecanismo moderno com suporte direto a
Unicode e fontes OpenType.

\begin{atencao}
  Não edite arquivos como \arquivo{.aux}, \arquivo{.log} ou \arquivo{.toc}.
  Eles são gerados durante a compilação. Seu conteúdo pertence aos arquivos
  \arquivo{.tex}.
\end{atencao}

\section{Comandos, argumentos e ambientes}

Um comando começa com barra invertida. Argumentos obrigatórios aparecem entre
chaves e opções aparecem entre colchetes:

\begin{lstlisting}[language=LaTeXBook]
\documentclass[11pt]{book}
\end{lstlisting}

Aqui, \arquivo{book} é obrigatório e \arquivo{11pt} é opcional. Ambientes têm
início e fim explícitos, como \comando{begin\{document\}} e
\comando{end\{document\}}. Se os nomes não coincidirem, a compilação falhará.

Alguns caracteres possuem significado especial: \verb|%| inicia comentário;
\verb|#| representa parâmetros de comandos; \verb|&| separa colunas. Para
imprimi-los no texto, geralmente usamos uma barra: \verb|\%|, \verb|\#| e
\verb|\&|.

\begin{pratica}
  Altere a classe para \arquivo{book}, acrescente a opção \arquivo{12pt} e
  escreva dois parágrafos separados por uma linha vazia. Compile novamente e
  observe o resultado.
\end{pratica}

\begin{resultado}
  Você já consegue identificar preâmbulo, corpo, comando, argumento e ambiente,
  além de transformar um arquivo-fonte em PDF.
\end{resultado}
